\section{Research Objectives and Questions}
This thesis addresses the inefficiencies in goal-oriented remote meetings through a structured approach that
integrates three complementary research objectives. The first objective (RQ1) investigates facilitating
strategies and challenges in remote meetings to establish a foundation based on real-world practices and pain
points. The second objective (RQ2) focuses on designing an AI-based moderation system that translates these
insights into technical requirements, ensuring the system is context-aware and goal-oriented. The third
objective (RQ3) evaluates user acceptance among knowledge workers.

The work bridges the gap between theoretical facilitation techniques and practical AI-driven solutions,
ensuring that the system is both technically robust and user-centric.
\subsection{RQ1: Facilitating strategies and challenges in goal-oriented remote meetings}
\begin{quote}
  \textit{Which facilitating strategies are suitable for the goal-oriented management of remote meetings with
    convergent objectives (such as decision-making, prioritization, problem-solving)? What specific challenges
  typically arise in this context?}
\end{quote}
The first research question aims to identify effective strategies for managing convergent remote meetings,
such as decision-making, prioritization, and problem-solving, while also analyzing the specific challenges
that arise. To achieve this, the thesis conducts expert interviews with certified facilitators, agile
coaches, and team leads to gather empirical insights. It also analyzes literature on facilitation techniques,
including timeboxing and role assignment, as well as the dynamics of remote collaboration. The findings are
compared with existing research to derive best practices and identify pain points, such as off-topic
discussions and unequal participation, which AI could potentially address.
\subsection{RQ2: Design requirements for an AI-based moderation system}
\begin{quote}
  \textit{How should an AI-based moderation system be designed that autonomously detects thematic deviations
    based on a predefined meeting goal and real-time transcript analysis and promotes goal orientation through
  constructive, context-sensitive interventions?}
\end{quote}
The second research question focuses on developing a proactive AI system that detects thematic deviations in
real time and intervenes contextually to maintain goal orientation. This objective is addressed by defining
functional requirements, such as real-time transcript analysis and adaptive interventions, and non-functional
requirements, including privacy, latency, and scalability. A system architecture is then designed with
components for transcription, NLP-based analysis, and intervention logic, incorporating mechanisms for
private nudges or group reminders. A prototype is implemented using NLP libraries like Hugging Face
Transformers, and its technical feasibility is evaluated.
\subsection{RQ3: User acceptance among knowledge workers}
\begin{quote}
  \textit{How do knowledge workers in remote teams evaluate the usefulness, usability, and psychological
  effects of such an AI moderation tool, and what factors influence its acceptance?}
\end{quote}
The third research question assesses how knowledge workers perceive the usefulness, usability, and
psychological impact of the AI moderation system, as well as the factors influencing its adoption. A
usability study is conducted with remote teams, using frameworks such as the System Usability Scale and the
Technology Acceptance Model. Quantitative data, including SUS scores and efficiency metrics, are collected
alongside qualitative feedback from interviews that explore perceptions of control, trust, and pressure. The
analysis also examines psychological effects, such as relief versus intrusion, and adoption barriers,
including resistance to automation and privacy concerns.
